\documentclass[11pt]{article}
\usepackage[margin=1in]{geometry}
\usepackage{amsmath,amssymb,amsthm}
\usepackage[hidelinks]{hyperref}
\newtheorem{theorem}{Theorem}
\newcommand{\C}{\mathbb C}
\newcommand{\norm}[1]{\left\lVert#1\right\rVert}
\title{A local power-series obstruction to a C*-algebraic Corona conjecture}
\author{Research proof and Lean formalization}
\date{October 27, 2026}
\begin{document}
\maketitle
\begin{abstract}
Under the local left-coefficient power-series definition used in the source,
a uniformly invertible affine function on the norm unit ball of $M_2(\C)$
need not have an analytic reciprocal in that class. A square-zero matrix
reduces the obstruction to two finite polynomial identities. This refutes
the source's C*-algebraic Corona assertion already for one datum. Ordinary
complex and operator-valued holomorphic Corona results are unaffected.
\end{abstract}

Krishna's Definition~2.6 in \cite{Krishna} calls a function analytic when,
near every center $a$, it has a norm-convergent expansion
\[
 f(z)=\sum_{n\geq0}b_n(z-a)^n
\]
with coefficients on the left. Conjecture~3.2 asserts that bounded functions
$f_1,\ldots,f_k$ in this class on the norm-open unit ball of a unital
C*-algebra, satisfying
\[
 \sum_j f_j(z)f_j(z)^*\geq\delta I,\qquad
 \sum_j f_j(z)^*f_j(z)\geq\delta I
\]
uniformly for some $\delta>0$, admit bounded functions $g_j,h_j$ of the
same class with $\sum_j f_jg_j=I=\sum_jh_jf_j$.
No commutativity assumption is imposed.

\begin{theorem}\label{main}
Let $\mathcal A=M_2(\C)$ with its operator norm, let
\[
 c=\begin{pmatrix}1/4&0\\0&0\end{pmatrix},\qquad f(z)=I+cz.
\]
On $\{z:\norm z<1\}$, the function $f$ is bounded and analytic in the
source's sense, and satisfies both lower bounds with $\delta=1/2$.
There is no function $g$ having a source power-series expansion even at
$0$ and satisfying $f(z)g(z)=I$ throughout that ball.
\end{theorem}

\begin{proof}
At any center $a$, the finite expansion $f(z)=f(a)+c(z-a)$ proves the
required analyticity. Since $\norm c=1/4$, we have $\norm{f(z)}\leq5/4$.
Set $T=cz$. The self-adjoint operator $T+T^*$ is bounded below by
$-2\norm T I$, and $TT^*,T^*T$ are positive. Thus
\[
 (I+T)(I+T)^*\geq(1-2\norm T)I\geq\tfrac12I,
 \qquad (I+T)^*(I+T)\geq\tfrac12I.
\]
Here $\norm T<1/4$ for every input in the ball. In fact the ordinary
reciprocal exists and is uniformly bounded by $4/3$, by the Neumann series.

Suppose $f(z)g(z)=I$ and $g(z)=\sum_{n\geq0}b_nz^n$ on some ball about
$0$. Substitution of $0$ gives $b_0=g(0)=I$. Put
\[
 N=\begin{pmatrix}1&1\\-1&-1\end{pmatrix}.
\]
Then
\begin{equation}\label{nil}
 N^2=0,\qquad
 cN=\begin{pmatrix}1/4&1/4\\0&0\end{pmatrix},\qquad
 (cN)^2=\begin{pmatrix}1/16&1/16\\0&0\end{pmatrix}\ne0.
\end{equation}
Choose a sufficiently small real $t>0$ so that both $tN$ and $-tN$ lie
in the expansion ball and the unit ball. Since all powers of either input
of order at least two vanish, norm convergence gives exactly
\[
 g(tN)=I+tb_1N,\qquad g(-tN)=I-tb_1N.
\]
Expanding the two right-inverse identities therefore yields
\[
 t(cN+b_1N)+t^2cNb_1N=0,\qquad
 -t(cN+b_1N)+t^2cNb_1N=0.
\]
Subtracting gives $b_1N=-cN$; adding gives $cNb_1N=0$.
Consequently $(cN)^2=0$, contradicting \eqref{nil}.
No differentiation or interchange of an infinite series is needed.
\end{proof}

\paragraph{Scope.}
Theorem~\ref{main} contradicts the conclusion of Conjecture~3.2 at $k=1$:
even its right solution cannot exist in the required class. The obstruction
concerns the prescribed local powers in a noncommutative algebra. The
ordinary Fr\'echet-holomorphic reciprocal exists; there is no contradiction
with scalar Carleson theory or with a different notion of noncommutative
analyticity. Rosenfeld\textquotesingle s matrix-algebra rigidity theorem already restricts
the relevant differentiable class to affine maps \cite[Theorem~2]{Rosenfeld}.
Restrictions on products are also discussed in \cite[Theorem~5.10]{Cook}.

\paragraph{Formal scope and status.}
The closed Lean endpoint \texttt{CStarCorona.source\_counterexample}
refutes the literal matrix-instance assertion, using the actual operator norm,
positive-semidefinite order and local norm-convergent left-power series.
The parent strict replay passes 25 declarations in four modules under
Lean~4.34.1 and mathlib revision
\[\texttt{d13f23b723b8a846827a245b89c10fc7d3f11612},\]
with only \texttt{propext}, \texttt{Classical.choice} and \texttt{Quot.sound}.
The ordinary Neumann reciprocal observation is classical background.

Discovery credit is zero: independent preliminary review confirms the
mathematics but rejects substantive eligibility as an immediate consequence
of classical rigidity. The later Lean proof and this note were not
independently reviewed. No new analytic machinery or worldwide priority is
claimed. The actual source was read in the author's Zenodo deposit;
direct Preprints PDF access remains unavailable.

\begin{thebibliography}{3}
\bibitem{Krishna} K.~Mahesh Krishna,
\emph{C*-algebraic Bieberbach, Robertson, Lebedev--Milin, Zalcman, Krzyz and
Corona Conjectures}, Preprints (2022), v1, Definition~2.6 and Conjecture~3.2,
pp.~15--16. \href{https://doi.org/10.20944/preprints202208.0140.v1}{doi:10.20944/preprints202208.0140.v1}.
Author deposit: \href{https://doi.org/10.5281/zenodo.6917925}{Zenodo 6917925}.
\bibitem{Cook} J.~S. Cook,
\emph{Introduction to $\mathcal A$-Calculus},
\href{https://arxiv.org/abs/1708.04135}{arXiv:1708.04135} (2017).
\bibitem{Rosenfeld} B. Rosenfeld,
\emph{Differentiable functions in associative and alternative algebras and
smooth surfaces in projective spaces over these algebras},
Publ. Inst. Math. (Beograd), N.S. \textbf{68}(82) (2000), 67--71.
\end{thebibliography}
\end{document}
